\documentclass[../main.tex]{subfiles}
\begin{document}
\chapter{Inner Product Spaces}
\begin{remark}[Note]
  Throughout this chapter we will be working with vector spaces over $\F = \R$ or $\F = \C$.
\end{remark}
\section{Definitions}
\begin{definition}[Real inner product]
  An \textit{inner product} $\varphi$ on an $\R$-vector space $V$ is a \textbf{positive definite symmetric} bilinear form on $V$.
\end{definition}
\begin{definition}[Complex inner product]
  An \textit{inner product} $\varphi$ on a $\C$-vector space $V$ is a \textbf{positive definite hermitian} bilinear form on $V$.
\end{definition}
\begin{definition}[Inner product space]
  For a vector space $V$ and inner product $\varphi$ on $V$, we call the pair $(V, \varphi)$ an \textit{inner product space}.
\end{definition}
\begin{remark}[Notation]
  We write $\langle u, v \rangle$ instead of $\varphi(u, v)$ when $\varphi$ is defines an inner product.
\end{remark}
\begin{example}
  \begin{enumerate}
    \item The standard dot product on $\R^{n}$ and $\C^{n}$ are real and complex inner products respectively.
    \item On the $\R$-vector space $V = C([0, 1])$, we can define:
      \[
        \langle f, g \rangle = \int_{0}^{1} f(t) g(t) \d{t}
      \]
      which is an inner product.
    \item More generally, given a \textit{weight function}:
      \[
        w: [0, 1] \to (0, \infty)
      \]
      we can define:
      \[
        \langle f, g \rangle_{w} = \int_{0}^{1} f(t) g(t) w(t) \d{t}
      \]
  \end{enumerate}
\end{example}
\begin{definition}[Length]
  For an inner product space $(V, \langle \argdot, \argdot \rangle)$, we define the \textit{length} of $v \in V$ as:
  \[
    \norm{v} = \sqrt{\langle v, v \rangle}
  \]
\end{definition}
\begin{definition}[Distance]
  For an inner product space $(V, \langle \argdot, \argdot \rangle)$, we define the \textit{distance} between $v \in V$ and $u \in V$ as $\norm{v - u}$.
\end{definition}
\begin{remark}[Remarks]
  \begin{enumerate}
    \item $\norm{v} = 0$ if and only if $v = \vec{0}$.
    \item $\norm{\argdot}: V \to \R_{\geq 0}$ uniquely determines $\langle \argdot, \argdot \rangle: V \times V \to \F$.
      This is because $\norm{v}^2 = \langle v, v \rangle$ and so we can use this to recover $\langle u, v \rangle$ for any $u, v \in V$ using the polarisation identity (\cref{polarisation} if $\F = \R$ and \cref{seqPolar} if $\F = \C$).
  \end{enumerate}
\end{remark}
\begin{lemma}[Cauchy-Schwarz]
  If $(V, \langle \argdot, \argdot \rangle)$ is an inner product space, then for any $v, w \in V$: \label{CS}
  \[
    |\langle v, w \rangle| \leq \norm{v} \norm{w}
  \]
\end{lemma}
\begin{remark}
  For the following proof, it is useful to note that symmetric inner products on $\R$-vector spaces are also hermitian as $\langle v, u \rangle \in \R \implies \langle u, v \rangle = \langle v, u \rangle = \overline{\langle v, u \rangle}$.
  This allows us to prove the above for both $\F = \C$ and $\F = \R$ by showing it for complex inner products.
\end{remark}
\begin{proof}
  If $v = \vec{0}$, then there is nothing to show as both sides are 0.
  Otherwise, for $t \in \F$:
  \begin{align*}
    0 &\leq \norm{tv + w}^2 \\
      &= \langle tv + w, tv + w \rangle \\
      &= |t|^2 \norm{v}^2 + t\langle v, w \rangle + \overline{t} \overline{\langle v, w \rangle} + \norm{w}^2
  \end{align*}
  If we set:
  \[
    t = - \frac{\overline{\langle v, w \rangle}}{\norm{v}^2} \in \F
  \]
  then:
  \[
    0 \leq \norm{w}^2 - \frac{|\langle v, w \rangle|^2}{\norm{v}^2} \implies |\langle v, w \rangle|^2 \leq \norm{v}^2 \norm{w}^2
  \]
  So taking square roots gives the desired result.
\end{proof}
\begin{remark}[Remarks]
  \begin{enumerate}
    \item We see that we have equality in Cauchy-Schwarz if and only if $\norm{tv + w} = 0$, i.e. if and only if $w = - tv$ which is if and only if the dimension of the span of $v$ and $w$ is 1.
    \item For $\F = \R$, using the bounds From Cauchy-Schwarz, we see that:
      \[
        \frac{\langle v, w \rangle}{\norm{v} \norm{w}} \in [-1, 1]
      \]
      and so we define the \textit{angle} $\theta$ between $v, w \in V$ as:
      \[
        \cos \theta = \frac{\langle v, w \rangle}{\norm{v}\norm{w}}
      \]
      which is consistent with our existing notion of angles in the case of $V = \R^2$.
  \end{enumerate}
\end{remark}
\begin{corollary}[Triangle inequality]
  If $(V, \langle \argdot, \argdot \rangle)$ is an inner product space, then for any $v, w \in V$:
  \[
    \norm{v + w} \leq \norm{v} + \norm{w}
  \]
\end{corollary}
\begin{proof}
  We can write:
  \begin{align*}
    \norm{v + w}^2 &= \langle v + w, v + w \rangle \\
                   &= \norm{v}^2 + \langle v, w \rangle + \overline{\langle v, w \rangle} + \norm{w}^2 \\
                   &\leq \norm{v}^2 + 2 \norm{v} \norm{w} + \norm{w}^2 \text{ by \cref{CS}} \\
                   &= (\norm{v} + \norm{w})^2
  \end{align*}
  So upon taking square roots, we obtain the desired result.
\end{proof}
\begin{remark}
  Since the triangle inequality holds, $d(v, w) = \norm{v - w}$ is a metric on $V$ as seen in IB Analysis II.
\end{remark}
\section{Orthogonality}
\subsection{Orthonormal Bases}
\begin{definition}[Orthogonal and Orthonormal]
  We say that a set $E \subseteq V$ is \textit{orthogonal} if for all $d, e \in E$:
  \[
    d \neq e \implies \langle d, e \rangle = 0
  \]
  Moreover, we say that $E \leq V$ is \textit{orthonormal} if it is orthogonal and $\norm{d} = 1$ for all $d \in V$.
\end{definition}
\begin{remark}
  If $V$ is a finite finite dimensional vector space and $B$ is a basis for $V$, then $B$ is orthonormal if and only if $[\langle \argdot, \argdot \rangle]_{B} = I_n$.
\end{remark}
\begin{lemma}
  If $E \subseteq V$ is orthogonal with $\vec{0} \notin E$, then: \label{vExpr}
  \begin{enumerate}
    \item If $v \in \langle E \rangle$, then:
      \[
        v = \finitesum_{e \in E} \frac{\langle v, e \rangle}{\langle e, e \rangle}e
      \]
    \item $E$ is linearly independent.
  \end{enumerate}
\end{lemma}
\begin{proof}[\textbf{(i)}]

  If $v \in \langle E \rangle$, then we can write:
  \[
    v = \finitesum_{e \in E} \lambda_e e
  \]
  for some $\lambda_e \in \F$.

  Thus for any $d \in E$:
  \begin{align*}
    \langle v, d \rangle &= \finitesum_{e \in E} \lambda_e \langle e, d \rangle \\
                         &= \lambda_d \langle d, d \rangle \text{ as $E$ is orthogonal}
  \end{align*}
  Since $d \neq \vec{0}$, we have:
  \[
    \lambda_d = \frac{\langle v, d \rangle}{\langle d, d \rangle}
  \]
  which gives the desired result.
\end{proof}
\begin{proof}[\textbf{ii}]
  Taking $v = \vec{0} \in \langle E \rangle$ in \textbf{(i)}, we see that:
  \[
    \lambda_d = \frac{\langle \vec{0}, d \rangle}{\langle d, d \rangle} = 0
  \]
  So:
  \[
    \finitesum_{e \in E} \frac{\langle v, e \rangle}{\langle e, e \rangle}e = \vec{0} \implies \lambda_e = 0\ \forall e \in E
  \]
  i.e. $E$ is linearly independent.
\end{proof}
\begin{corollary}[Parseval's Identity]
  If $E \subseteq V$ is orthonormal and $v, w \in \langle E \rangle$, then:
  \[
    \langle v, w \rangle = \finitesum_{e \in E} \langle v, e \rangle \overline{\langle w, e \rangle}
  \]
\end{corollary}
\begin{proof}
  Since $E$ is orthonormal, $\vec{0} \notin E$ and $\norm{\vec{0}} \neq 1$ so we can apply \cref{vExpr}.
  By orthonormality $\langle e, e \rangle = 1$ for all $e \in E$ and so we can write any $v, w \in V$ as:
  \[
    v = \finitesum_{e \in E} \langle v, e \rangle e \text{ and } w = \finitesum_{d \in E} \langle w, d \rangle d
  \]
  Thus:
  \begin{align*}
    \langle v, w \rangle &= \left\langle \finitesum_{e \in E} \langle v, e \rangle e, \finitesum_{d \in E} \langle w, d \rangle d \right\rangle \\
                         &= \finitesum_{d, e \in E} \langle v, e \rangle \overline{\langle w, d \rangle} \langle e, d \rangle \text{ by anti-linearity and linearity} \\
                         &= \finitesum_{e \in E} \langle v, e \rangle \overline{\langle w, e \rangle} \text{ as $\langle e, d \rangle = \delta_{e d}$ by orthonormality}
  \end{align*}
  as required.
\end{proof}
\begin{theorem}[Gram-Schmidt Process]
  If $(V, \langle \argdot, \argdot \rangle)$ is an inner product space and $\{v_1, \ldots, v_n\} \subseteq V$ is linearly independent, then there is an orthonormal set $\{e_1, \ldots, e_n\} \subseteq V$ such that: \cref{GS}
  \[
    \langle e_1, \ldots, e_k \rangle = \langle v_1, \ldots, v_k \rangle
  \]
  for each $k \in \{1, \ldots, n\}$.
\end{theorem}
\begin{proof}
  We proceed by induction on $k$.
  In the case of $k = 1$, we set $e_1 = \frac{1}{\norm{v_1}} v_1$, clearly $\langle e_1 \rangle = \langle v_1 \rangle$ and $\{e_1\}$ is orthonormal.

  Suppose we have constructed $e_1, \ldots, e_k$ as specified for some $k \leq n - 1$.
  Define:
  \[
    \widetilde{e}_{k + 1} = v_{k + 1} - \sum_{i = 1}^{k} \langle v_{k + 1}, e_i \rangle e_i
  \]
  We then observe that $\{v_1, \ldots, v_k, v_{k + 1}\}$ is linearly independent as it is a subset of a linearly independent set and so $v_{k + 1} \notin \langle v_1, \ldots, v_k \rangle = \langle e_1, \ldots, e_k \rangle$.
  This tells us that $\widetilde{e}_{k + 1} \neq \vec{0}$.

  We also see that for any $j \in \{1, \ldots, k\}$:
  \begin{align*}
    \langle \widetilde{e}_{k + 1}, e_j \rangle &= \langle v_{k + 1}, e_j \rangle - \langle v_{k + 1}, e_j \rangle \\
                                               &= 0
  \end{align*}
  since $\{e_1, \ldots, e_k\}$ is an orthonormal set, thus $\{e_1, \ldots, e_k, \widetilde{e}_{k + 1}\}$ is orthogonal.

  Finally, as $\widetilde{e}_{k + 1}$ is a linear combination of $v_{k + 1}$ and an element in $\langle e_1, \ldots, e_k \rangle$, we have that:
  \[
    \langle e_1, \ldots, e_k, \widetilde{e}_{k + 1} \rangle = \langle v_1, \ldots, v_{k + 1} \rangle
  \]

  The last thing we need to do is normalise $\widetilde{e}_{k + 1}$ which we can do as $\widetilde{e}_{k + 1} \neq \vec{0}$, i.e. we let:
  \[
    e_{k + 1} = \frac{\widetilde{e}_{k + 1}}{\norm{\widetilde{e}_{k + 1}}}
  \]
  Thus the set $\{e_1, \ldots, e_k, e_{k + 1}\}$ is orthogonal with $\langle e_1, \ldots, e_{k + 1} \rangle = \langle v_1, \ldots, v_{k + 1} \rangle$, as required.
\end{proof}
\begin{remark}
  Suppose $\{v_1, \ldots, v_k\}$ is already orthonormal for some $k \leq n$, then we see that $e_1 = v_1 \implies e_2 = v_2 \implies \cdots \implies e_k = v_k$, i.e. $e_i = v_i$ for $i \in \{1, \ldots, k\}$.
\end{remark}
\begin{corollary}
  If $(V, \langle \argdot, \argdot \rangle)$ is an inner product space on a finite vector space $V$, then any orthonormal set in $V$ can be extended to an orthonormal basis for $V$. \label{orthoBasis}
\end{corollary}
\begin{proof}
  Given an orthonormal set $\{v_1, \ldots, v_k\}$, we use \cref{basisExtension} to extend it to a basis for $V$ and then apply \cref{GS} to turn it into an orthonormal set with the same span as the original set.
  Moreover, as the original set is already orthonormal, $e_i = v_i$ for $i \in \{1, \ldots, k\}$ and thus we have extended our original orthonormal set to an orthonormal basis.
\end{proof}
\begin{remark}
  In particular, $V \neq \{\vec{0}\}$ always has an orthonormal basis because we can take any $v \in V \setminus \{\vec{0}\}$, normalise it to length 1 so that the set $\{v\}$ is orthonormal and then use the above to extend it to an orthonormal basis for $V$.
\end{remark}
\subsection{Orthogonal Sums and Projections}
\begin{definition}[Orthogonal direct sum]
For an inner product space $(V, \langle \argdot, \argdot \rangle)$ and $U, W \leq V$, we say that $V$ is the \textit{orthogonal direct sum} of $U$ and $W$ if:
  \begin{enumerate}
    \item $V = U + W$
    \item $\langle u, w \rangle = 0$ for all $u \in U$ and $w \in W$
  \end{enumerate}
  We then write $V = U \obot W$.
\end{definition}
\begin{remark}
  Suppose $V = U \obot W$ and $U \cap W \neq \{\vec{0}\}$.
  Then by \textbf{(ii)} there exists $v \in U \cap W \setminus \{\vec{0}\}$ such that $\langle v, v \rangle = 0 \implies v = \vec{0}$, which is a contradiction.
  Thus:
  \[
    V = U \obot W \implies U \cap W = \{\vec{0}\} \implies V = U \oplus W
  \]
\end{remark}
\begin{lemma}
  If $V$ is finite dimensional and $W \leq V$, then we can write $V = W \obot W^{\perp}$.
\end{lemma}
\begin{proof}
  By definition of $W^{\perp}$, $\langle u, w \rangle = 0$ for all $u \in W$ and $v \in W^{\perp}$, so \textbf{(ii)} is satisfied.

  Since $W$ is finite dimensional, by the remark of \cref{orthoBasis}, there exists an orthonormal basis $E_W$ for $W$.
  Moreover, by \cref{orthoBasis}, we can extend this to an orthonormal basis $E_V$ for $V$, then by construction:
  \[
    V = W + \langle E_V \setminus E_W \rangle
  \]
  Since $E_V$ is orthogonal, all the vectors in $E_V \setminus E_W$ are orthogonal to those in $E_W$ and thus everything in $\langle E_v \setminus E_W \rangle$ is orthogonal to everything in $\langle E_W \rangle = W$, so $\langle E_V \setminus E_W \rangle \subseteq W^{\perp}$.
  Thus:
  \[
    V = W + \langle E_V \setminus E_W \rangle \subseteq W + W^{\perp} \subseteq V
  \]
  which tells us that $V = W + W^{\perp}$, thus \textbf{(i)} is satisfied also.

  Both conditions are satisfied so $V = W \obot W^{\perp}$.
\end{proof}
\begin{corollary}
  If $V$ is finite dimensional then for any $W \leq V$:
  \[
    \dim(V) = \dim(W) + \dim(W^{\perp})
  \]
\end{corollary}
\begin{proof}
  As discussed $V = W \obot W^{\perp} \implies V = W \oplus W^{\perp}$.
  So by \cref{externalDim}:
  \[
    \dim(V) = \dim(W \oplus W^{\perp}) = \dim(W) + \dim(W^{\perp})
  \]
  as required.
\end{proof}
\begin{definition}[Orthogonal projection]
  Let $(V, \langle \argdot, \argdot \rangle)$ be an inner product space with $W \leq V$.
  Since $V = W \oplus W^{\perp}$, for any $v \in V$ there is a unique pair of elements $w \in W$ and $u \in W^{\perp}$ such that $v = w + u$.
  We use this to define the \textit{projection of $V$ onto $W$} as:
  \[
    \pi_W: V \to V,\ v \mapsto w
  \]
\end{definition}
\begin{remark}
  \begin{enumerate}
    \item Suppose $v_1 = w_1 + u_1$ and $v_2 = w_2 + u_2$ as above, then for any $\lambda \in \F$:
      \begin{align*}
        \pi_W(v_1 + \lambda v_2) &= \pi_W((w_1 + \lambda w_2) + (u_1 + \lambda u_2)) \\
                                 &= w_1 + \lambda w_2 \\
                                 &= \pi_W(v_1) + \lambda \pi_W(v_2)
      \end{align*}
      So $\pi_W$ is linear.
    \item $\pi_W$ only outputs elements of $W$ and if $w \in W$, then $\pi_W(w) = w$, thus $\im \pi_W = W$.
    \item Write $v = w + u$ and suppose $v \in \ker \pi_W$, then $\pi_W(v) = w = \vec{0}$, so $v = u \in W^{\perp}$.
      Similarly $v \in W^{\perp} \implies v \in \ker \pi_W$.
      Thus $\ker \pi_W = W^{\perp}$.
    \item $\at{\pi_{W}}{W} = \id_W$, i.e. $(\pi_W)^2 = \pi_W$.
    \item For any $v \in V$ write $v = w + u$, then:
      \begin{align*}
        (\id_V - \pi_W)[v] &= w + u - w \\
                           &= u \\
                           &= \pi_{W^{\perp}}(v)
      \end{align*}
      So $\id_V - \pi_W = \pi_{W^{\perp}}$.
  \end{enumerate}
\end{remark}
\begin{proposition}[Closest point projection]
  For all $v \in V$ and $w \in W \leq V$:
  \[
    \norm{v - \pi_W(v)} \leq \norm{v - w}
  \]
  with equality if and only if $w = \pi_W(v)$.
\end{proposition}
\begin{remark}[Intuition]
  For any point $v \in V$ and a subspace $W \leq V$, the closest point in $W$ to $v$ is its projection onto $W$, i.e. $\pi_W(v)$.
\end{remark}
\begin{proof}
  We see that:
  \begin{align*}
    \norm{v - w}^2 &= \norm{(v - \pi_{W}(v)) - (\pi_W(v) - w)}^2
  \end{align*}
  Let $u_1 = v - \pi_W(v) \in W^{\perp}$ and $u_2 = \pi_W(v) - w \in W$, then:
  \begin{align*}
    \norm{v - w}^2 &= \langle u_1 - u_2, u_1 - u_2 \rangle \\
                   &= \langle u_1, u_1 \rangle - 2\cancelto{0}{\langle u_1, u_2 \rangle} + \langle u_2, u_2 \rangle \\
                   &= \norm{u_1}^2 + \norm{u_2}^2 \\
                   &= \norm{v - \pi_W(v)}^2 + \norm{\pi_{W}(v) - w}^2 \\
                   &\geq \norm{v - \pi_W(v)}^2
  \end{align*}
  with equality if and only if $\norm{\pi_W(v) - w}^2 = 0 \iff w = \pi_{W}(v)$.
\end{proof}
\begin{definition}[Inner product on the external direct sum]
  If $(V_1, \langle \argdot, \argdot \rangle_1)$ and $(V_2, \langle \argdot, \argdot \rangle_2)$ are inner product spaces, then:
  \[
    \langle (v_1, v_2), (w_1, w_2) \rangle = \langle v_1, w_1 \rangle_1 + \langle v_2, w_2 \rangle_2
  \]
  defines an inner product on the external direct sum $V_1 \oplus V_2$.
\end{definition}
\begin{proof}
  We can check that this satisfies the criterion required of an inner product.
  \begin{enumerate}
    \item Linearity and anti-linearity are inherited from $\langle \argdot, \argdot \rangle_1$ and $\langle \argdot, \argdot \rangle_2$.
    \item Hermitianity is inherited from $\langle \argdot, \argdot \rangle_1$ and $\langle \argdot, \argdot \rangle_2$.
    \item Suppose $(v_1, v_2) \neq \vec{0}$ then:
      \[
        \langle (v_1, v_2), (v_1, v_2) \rangle = \langle v_1, v_1 \rangle_1 + \langle v_2, v_2 \rangle_2 > 0
      \]
      Moreover:
      \begin{align*}
        \langle (v_1, v_2), (v_1, v_2) \rangle = 0 &\iff \langle v_1, v_1 \rangle_1 = 0 \text{ and } \langle v_2, v_2 \rangle_2 = 0 \\
                                                   &\iff v_1 = \vec{0} \text{ and } v_2 = \vec{0} \\
                                                   &\iff (v_1, v_2) = \vec{0}
      \end{align*}
      So $\langle \argdot, \argdot \rangle$ is positive definite.
  \end{enumerate}
\end{proof}
\section{Adjoint Maps}
\begin{remark}[Note]
  Henceforth, assume that $(V, \langle \argdot, \argdot \rangle_{V})$ and $(W, \langle \argdot, \argdot \rangle_{W})$ are finite dimensional inner product spaces.
\end{remark}
\begin{theorem}
  If $\alpha \in \mathcal{L}(V, W)$, and $B$ and $C$ are orthonormal bases for $V$ and $W$ respectively then: \label{adjUnique}
  \begin{enumerate}
    \item There exists a unique linear map $\alpha^{*} : W \to V$ such that:
      \begin{equation}
        \langle \alpha(v), w \rangle_{W} = \langle v, \alpha^{*}(w) \rangle_{V}\ \forall v \in V, w \in W \label{adjProp}
      \end{equation}
    \item Moreover:
      \[
        [\alpha^{*}]^{C}_{B} = (\overline{[\alpha]^{B}_{C}})^{\trans}
      \]
  \end{enumerate}
\end{theorem}
\begin{definition}[Adjoint map]
  We call $\alpha^{*}$ the \textit{adjoint map} of $\alpha$.
\end{definition}
\begin{remark}[Warning]
  You might notice that this is the same notation as the dual map $\alpha^{*}: W^{*} \to V^{*}$, there is a reason for this that we will see later.
\end{remark}
\begin{remark}[Recap]
  Recall that if $\dim V = n$ and $B$ is an orthonormal basis for $V$, then:
  \[
    [\langle \argdot, \argdot \rangle]_{B} = I_n
  \]
\end{remark}
\begin{proof}[\textbf{i} -- \cref{adjUnique}]
  Write $B = \{b_1, \ldots, b_n\}$, $C = \{c_1, \ldots, c_n\}$, and set $A = [\alpha]^{B}_{C}$.

  \textbf{Existence}\par
  Define $\alpha^{*}: W \to V$ by:
  \[
    \alpha^{*}(c_j) = \sum_{i = 1}^{n} \overline{a_{j i}} b_i
  \]
  so that:
  \[
    [\alpha^{*}]^{C}_{B} = \overline{A}^{\trans}
  \]
  Then for any $v \in V$ and $w \in W$:
  \begin{align*}
    \langle \alpha(v), w \rangle_{W} &= ([\alpha(v)]_C)^{\trans} [\langle \argdot, \argdot \rangle_W]_C \overline{[w]_{C}} \text{ by \cref{seqMat}} \\
                                     &= (A[v]_{B})^{\trans}\overline{[w]_{C}} \text{ by \cref{basisMapProp} \textbf{(i)}} \\
                                     &= [v]^{\trans}_B (A^{\trans} \overline{[w]_C}) \\
                                     &= [v]^{\trans}_B \overline{(\overline{A}^{\trans} [w]_C)} \\
                                     &= [v]^{\trans}_B \overline{([\alpha^{*}]^{C}_{B} [w]_C)} \text{ by construction} \\
                                     &= [v]^{\trans}_B \overline{[\alpha^{*}(w)]_B} \text{ by \cref{basisMapProp} \textbf{(i)}} \\
                                     &= [v]^{\trans}_B [\langle \argdot, \argdot \rangle_V]_B \overline{[\alpha^{*}(w)]_B} \\
                                     &= \langle v, \alpha^{*}(w) \rangle_{V}
  \end{align*}
  Thus such an adjoint map $\alpha^{*}$ exists.

  \textbf{Uniqueness}\par
  Suppose $\alpha': W \to V$ also satisfies \cref{adjProp} and let $A' = [\alpha']^{C}_{B}$.
  Similarly to how we showed existence:
  \begin{align*}
    \langle \alpha(b_i), c_j \rangle_W &= [\alpha(b_i)]^{\trans}_C [\langle \argdot, \argdot \rangle_W]_C \overline{[c_j]_C} \text{ by \cref{seqMat}} \\
                                       &= ([\alpha]^{B}_{C}[b_i]_B) I_n e_j \\
                                       &= (Ae_i)^{\trans} e_j \\
                                       &= e^{\trans}_i A^{\trans} e_j \\
                                       &= a_{j i}
  \end{align*}
  and also:
  \begin{align*}
    \langle b_i, \alpha'(c_j) \rangle_V &= [b_i]^{\trans}_B [\langle \argdot, \argdot \rangle_V]_B \overline{[\alpha'(c_j)]_B} \text{ by \cref{seqMat}} \\
                                      &= e^{\trans}_i I_n \overline{[\alpha']^{C}_{B}[c_j]_C} \\
                                      &= e^{\trans}_i \overline{[\alpha']^{C}_{B} e_j} \\
                                      &= e^{\trans}_i \overline{[\alpha']^{C}_{B}} e_j \\
                                      &= e^{\trans}_i \overline{A'} e_j \\
                                      &= \overline{a_{i j}'} \\
  \end{align*}
  As $\alpha'$ satisfies \cref{adjProp}:
  \[
    a_{j i} = \langle \alpha(b_i), c_j \rangle_W = \langle b_i, \alpha'(c_j) \rangle_V = \overline{a_{i j}'}
  \]
  Thus:
  \[
    [\alpha']^{C}_{B} = A' = \overline{A}^{\trans} = [\alpha^{*}]^{C}_{B}
  \]
  so $\alpha' = \alpha^{*}$, i.e. $\alpha$ is unique.
\end{proof}
\begin{proof}[\textbf{ii} -- \cref{adjUnique}]
  In the proof for \textbf{(i)}, we constructed $\alpha^{*}$ to have this property, and thus as it is unique by \textbf{(i)}, the adjoint map $\alpha^{*}$ always has this property.
\end{proof}
\begin{remark}[Notation]
  We have been using the notation $\overline{A}^{\trans}$ regularly for the conjugate transpose (i.e. the Hermitian conjugate), for brevity, we introduce the shorthand $A^{\dag} = \overline{A}^{\trans}$.

  In this notation, we have:
  \[
    [\alpha^{*}]^{B}_{C} = ([\alpha]^{B}_{C})^{\dag}
  \]
\end{remark}
\begin{example}
  Recalling that $\F = \R$ or $\F = \C$, for $v \in \F^{n}, w \in \F^{m}$, and $A \in \mathcal{M}_{m \times n}(\F)$, then, considering the standard dot product on $\R$ or $\C$ respectively:
  \begin{align*}
    (Av) \cdot w &= (Av)^{\trans} \overline{w} \\
                 &= v^{\trans} (A^{\trans} \overline{w}) \\
                 &= v^{\trans} \overline{(A^{\dag} w)} \\
                 &= v \cdot (A^{\dag} w)
  \end{align*}
  So the adjoint of $A$ is its Hermitian conjugate $A^{\dag}$.
\end{example}
\begin{remark}[Recap]
  Recall from \cref{linRightMap} that if $(V, \langle \argdot, \argdot \rangle)$ is a \textbf{real} inner product space and $\varphi(v, w) = \langle v, w \rangle$, then the map $\varphi_{R, V}: V \to V^{*}$ is linear where for any $v \in V$:
  \[
    \varphi_{R, V}(v)[u] = \langle u, v \rangle \text{ for } u \in V
  \]
  Note that we require $\F = \R$ as \cref{linRightMap} is valid only for bilinear forms.
  Inner products on $\C$ are sesquilinear forms so do not have to be bilinear.
\end{remark}
\begin{theorem}[Riesz representation]
  If $(V, \langle \argdot, \argdot \rangle)$ is a finite dimensional \textbf{real} inner product space, then $\varphi_{R, V}: V \to V^{*}$ is an isomorphism.
\end{theorem}
\begin{proof}
  We saw in \cref{dualSize} that $\dim(V^{*}) = \dim(V)$ and so by \cref{linPidg}, it suffices to show that $\varphi_{R, V}$ is linear and injective.

  \textbf{Linearity -} See the above remark.

  \textbf{Injectivity -} If $v \in \ker \varphi_{R,V}$, then as $v \in V$:
  \[
    0 = \varphi_{R, V}(v)[v] = \langle v, v \rangle
  \]
  Since $\langle \argdot, \argdot \rangle$ is positive definite, we must have $v = \vec{0}$, i.e. $\ker \varphi_{R, V}$ is trivial, so $\varphi_{R, V}$ is injective.
\end{proof}
\end{document}
